% !TEX root = 00_main_aaai_2027.tex
\section{Introduction}
A scalar agreement statistic does not determine whether a judge predicts how human ratings vary across inputs.
G-Eval evaluates text judges using correlation with human judgments~\citep{liu2023geval}.
MT-Bench and Chatbot Arena also study agreement with human preferences and document judge biases~\citep{zheng2023judging}.
These comparisons answer useful questions, but do not by themselves identify a conditional rating distribution.

We take the distribution as the target when the evaluation requires probabilities for each possible rating.
Proper scoring rules provide an established framework for this target~\citep{gneiting2007strictly}.
Our contribution is a controlled identification study that makes the information discarded by scalar criteria explicit.
We do not claim a new proper scoring rule.

The study separates three questions: population identification, selection with sparse ratings, and agreement of scalar means.
It compares the same candidates and splits under probability-loss and correlation selection.
A separate permutation control shows why matching aggregate histograms cannot establish correct item conditioning.

We make the following contributions:
\begin{enumerate}
\item We construct conditional distributions that share their means but differ under proper probability losses.
\item We evaluate selection using one, two and three synthetic ratings per item, reporting every tied minimizer and common test metric.
\item We preserve historical null findings and identify the limits of transferring this synthetic result to human-rated judges.
\end{enumerate}
